% !TeX root = ../../../manual.tex
\chapter{Bibliography}

\section{biblatex + biber Setup}

OmniLaTeX configures \texttt{biblatex} with \texttt{biber} as the backend.
The default style is \texttt{ext-authoryear} with several OmniLaTeX-specific
customisations:

\begin{verbatim}
\usepackage[
    backend      = biber,
    style        = ext-authoryear,
    sorting      = nyt,
    maxbibnames  = 99,
    maxcitenames = 2,
    giveninits   = true,
    uniquename   = init,
    uniquelist   = false,
    dashed       = false,
    isbn         = false,
    url          = false,
    doi          = false,
    eprint       = false,
    backref      = true,
    backrefstyle = none,
]{biblatex}
\end{verbatim}

Key OmniLaTeX customisations:

\begin{itemize}
    \item \textbf{introcite=label} --- prints a numeric label at the first
          citation in each chapter.
    \item \textbf{Family-given name format} --- author names are rendered as
          \textsc{Family} Given.
    \item \textbf{Semicolon delimiter} --- multiple authors in a single
          citation are separated by semicolons.
    \item \textbf{Small caps family names} --- family names appear in small
          capitals for visual distinction.
\end{itemize}

\section{Adding Bibliography}

Declare bibliography resource files with \verb|\addbibresource|:

\begin{verbatim}
\addbibresource{references.bib}
\addbibresource{supplementary.bib}
\end{verbatim}

Multiple \texttt{.bib} files can be loaded. All entries from all files are
available for citation. The \texttt{.bib} files are parsed by \texttt{biber}
at compile time.

Print the bibliography at the end of the document:

\begin{verbatim}
\printbibliography
\end{verbatim}

\section{Citation Commands}

OmniLaTeX supports the full set of \texttt{biblatex} citation commands:

\begin{longtblr}[
    caption = {Citation commands},
    label = {tab:citation-commands},
]{colspec = {l l}, rowhead = 1}
    \toprule
    Command              & Output \\
    \midrule
    \verb|\cite{key}|        & Numeric label: [1] \\
    \verb|\textcite{key}|    & Author (year) inline \\
    \verb|\parencite{key}|   & Parenthesised: (Author, year) \\
    \verb|\footcite{key}|    & Footnote citation \\
    \verb|\fullcite{key}|    & Full bibliographic entry inline \\
    \verb|\autocite{key}|    & Context-sensitive (parencite or footcite) \\
    \verb|\citeauthor{key}|  & Author name only \\
    \verb|\citeyear{key}|    & Year only \\
    \verb|\parencites{key1}{text1}| & Multiple parenthesised citations \\
    \verb|\textcites{key1}{text1}|  & Multiple textual citations \\
    \verb|\cites{key1}{text1}|      & Multiple citations with notes \\
    \verb|\footcites{key1}{text1}|  & Multiple footnote citations \\
    \bottomrule
\end{longtblr}

\section{Citation Groups}

Group multiple citations with per-citation notes using the starred variants:

\begin{verbatim}
\parencites{key1}[see][p.~42]{key2}[cf.]{key3}

\cites{key1}{see also \textit{key2}}{key3}

\textcites{key1}{as discussed in key2}{key3}
\end{verbatim}

Each key may be followed by optional pre- and post-note arguments in square
brackets: \verb|\parencites[pre1][post1]{key1}[pre2][post2]{key2}|.

\section{Per-Chapter Bibliography}

Two mechanisms exist for splitting the bibliography by chapter:

\begin{description}
    \item[\texttt{refsection}] --- Creates independent bibliography
        sections. Each \texttt{refsection} environment resets the citation
        counter and produces its own bibliography. Use
        \verb|\printbibliography[heading=subbibliography]| inside each
        section.

\begin{verbatim}
\begin{refsection}
    \chapter{Thermodynamics}
    % ... citations ...
    \printbibliography[heading=subbibliography]
\end{refsection}
\end{verbatim}

    \item[\texttt{refsegment}] --- Segments share a single bibliography but
        group entries by segment. Use \verb|\printbibliography[
        segment=\therefsegment]| to list only the current segment's entries.
\end{description}

\texttt{refsegment} is simpler to set up; \texttt{refsection} provides
stronger isolation between chapters.

\section{Back References}

The \texttt{backref} option adds page references to each bibliography entry
showing where it was cited:

\begin{verbatim}
\usepackage[backref=true, backrefstyle=none]{biblatex}
\end{verbatim}

\begin{longtblr}[
    caption = {Back reference styles},
    label = {tab:backref-styles},
]{colspec = {l l}, rowhead = 1}
    \toprule
    Style     & Description \\
    \midrule
    none      & Cited on pages 1, 3, 7 \\
    chapter   & Cited on chapters 1, 2 \\
    section   & Cited on sections 1.1, 2.3 \\
    \bottomrule
\end{longtblr}

\section{Further Reading Section}

A ``Further Reading'' section lists uncited references using
\texttt{category} filtering:

\begin{verbatim}
\addbibresource{cited.bib}
\addbibresource{uncited.bib}

% In cited.bib:  @book{key1, ...}
% In uncited.bib: @book{key2, category={notcited}, ...}

\printbibliography[title={References}, notcategory=notcited]
\printbibliography[title={Further Reading}, category=notcited]
\end{verbatim}

Alternatively, use \verb|\nocite{*}| to include all entries, then split by
category. Entries without a \texttt{category} field default to
\texttt{cited}.
